\indent

\textbf{This section examines the key features under my responsibility in the GMES School system}, analyzing user role interactions and business logic.
% COMMON
\subsubsection{Common Elements and Flows for All Use Cases}\label{sec:common_usecase}
\par\indent\indent\parbox{0.9\textwidth}{
    This section establishes the foundational elements that are consistent across all use cases documented in \textbf{\autoref{sec:use_case_specification} \nameref{sec:use_case_specification}}
}\\[0.25cm]
\indent\textbf{2.4.1.1 Standard Pre-conditions}\\
\indent\indent User must be logged into the system with valid credentials.\\
\indent\indent User must have appropriate role-based permissions for the specific function.\\[0.2cm]
\indent\textbf{2.4.1.2 Standard Request Data}\\
\indent\indent User session data (token, context).\\[0.2cm]
\indent\textbf{2.4.1.3 Standard Alternative Flow}
\vspace{0.2cm}
\par\indent\indent\parbox{0.9\textwidth}{\textbf{Network/server error}: If network/server errors occur, the system displays an alert about connectivity issues, notifying users when features are unavailable. This allows users to retry or report the problem to support staff.}

% IMPORT EMPLOYEE DATA
\subsubsection{Personnel Management: Import Employee Data}
\indent\indent\textbf{2.4.2.1 Brief description}
\vspace{0.5em}
\par\indent\indent\parbox{0.9\textwidth}{
    This use case demonstrates how school administrators and staff import employee information using Excel files through the system's web application.
}\\

\vspace{0.5em}
\indent\textbf{2.4.2.2 Flow of event}
\vspace{0.5em}
\par\indent\indent\textbf{A. Basic flow}
\vspace{-0.5em}
\renewcommand{\arraystretch}{1.5}
\begin{longtable}{|>{\raggedright\arraybackslash}p{4.5cm}|
                  >{\raggedright\arraybackslash}p{6cm}|
                  >{\raggedright\arraybackslash}p{4.5cm}|}                  
\caption{Import employee data via Excel basic flow.}
\label{tab:employee_management_basic_flow1} \\
\hline
\centering\textbf{User action} & \centering\textbf{System action} & \centering\arraybackslash\textbf{Data} \\
\hline
\endfirsthead
\endhead

1. User navigates to the Employee profiles tab. & 2. The system retrieves the organizations that are managed by the user.& \\
    \hline
    3. User selects a department and clicks the "Upload employee file" button. &4. System prepares to receive the Excel file.&\\
    \hline
    5. User uploads the Excel file.& 6. System validates the file format and checks for required columns.
    &\par{Request:
        \vspace{-0.8em}
        \begin{itemize}
            \item Excel file
        \end{itemize}
    } \\
    \hline
    &7. The system reads the file and checks each employee code to determine whether the employee is new or existing.&\\
    \hline
    &8. System verifies all data fields in each row.&\\
    \hline
    &9. If the employee is new, the system adds a new record to the database, otherwise, the system updates the existing employee information.&\\
    \hline
    &10. The system reloads and displays the latest employee list for the chosen department.&\\
    \hline
\end{longtable}

\par\indent\indent\textbf{B. Alternative flow}
\vspace{-0.5em}
\renewcommand{\arraystretch}{1.5}
\begin{longtable}{|>{\raggedright\arraybackslash}p{4.5cm}|
                  >{\raggedright\arraybackslash}p{6cm}|
                  >{\raggedright\arraybackslash}p{4.5cm}|}                  
\caption{Import employee data via Excel alternative flow.}
\label{tab:employee_management_import_data_alternative_flow} \\
\hline
\centering\textbf{User action} & \centering\textbf{System action} & \centering\arraybackslash\textbf{Data} \\
\hline
\endfirsthead
\endhead

5. User uploads the Excel file.& 6a. System detects invalid file format (not .xlsx/.xls).
&\par{Request:
    \vspace{-0.8em}
    \begin{itemize}
        \item Excel file
        \vspace{-1.4em}
    \end{itemize}
}\\

\hline
&6b. The system checks for missing required fields such as employee code, name, ID number, work office, etc.&
\par{Request:
    \vspace{-0.8em}
    \begin{itemize}
        \item Excel file
    \end{itemize}
}\\
\hline
&6c. For valid file, system detects a duplication in ID number.&\\
\hline
&6d. System identifies invalid data in a row such as wrong date format.& \\
\hline
&6e. System detects reference data mismatch such as invalid department code.&\\
\hline
&7. The system updates the valid data that it read before encountering the error, the rest of the file will be ignored. &\\
\hline
&10. The system displays an error message.&\\
\hline
\end{longtable}
\newpage
\indent\textbf{2.4.2.3 Notes}\\
\vspace{-0.7em}
\par\indent\indent\parbox{0.9\textwidth}{
    The uploaded Excel file must follow the system's predefined template. User can download the template from the web application of the system.
}\\[0.2cm]
\indent\textbf{2.4.2.4 Pre-conditions}\\
\indent\indent The uploaded file must be Excel file (.xlsx/.xls).\\[0.2cm]
\indent\textbf{2.4.2.5 Post-conditions}\\
\vspace{-0.7em}
\par\indent\indent\parbox{0.9\textwidth}{
    On success, the system updates all employee data from the file to the database and refreshes the data. If an error occurs, the system updates only the valid data before it encounters the error and displays an error message.
}\\

% STUDENT ATTENDANCE MANAGEMENT
\subsubsection{Attendance Management: Mark Student Attendance}
\indent\indent\textbf{2.4.3.1 Brief description}
\vspace{0.5em}
\par\indent\indent\parbox{0.9\textwidth}{
    This use case demonstrates how teacher and bus managers mark student attendance through the mobile app, and how the system handles the request.
}\\

\vspace{0.5em}
\indent\textbf{2.4.3.2 Flow of event}
\vspace{0.5em}
\par\indent\indent\textbf{A. Basic flow}
\vspace{1em}
\par\indent\indent\indent\parbox{0.87\textwidth}{\textbf{Mark student attendance in class}: This feature enables teachers to record student attendance when they in class.}
\vspace{-1em}
\renewcommand{\arraystretch}{1.5}
\begin{longtable}{|>{\raggedright\arraybackslash}p{4.5cm}|
                  >{\raggedright\arraybackslash}p{6cm}|
                  >{\raggedright\arraybackslash}p{4.5cm}|}                  
\caption{Mark student attendance basic flow.}
\label{tab:student_attendance_mark_presence_class_basic_flow}\\
\hline
\centering\textbf{User action} & \centering\textbf{System action} & \centering\arraybackslash\textbf{Data} \\
\hline
\endfirsthead
\endhead

1. User selects a class and navigates to Attendance function. & 2. The system retrieves and display a list of students with their current status (checked/unchecked).& 
    \par{Request:
        \vspace{-0.8em}
        \begin{itemize}
            \item Id of the selected class
        \end{itemize}
    } 
    \vspace{-0.8em}
    \par{Response: 
        \vspace{-0.8em}
        \begin{itemize}
            \item A list of student.
            \vspace{-1.5em}
        \end{itemize}
    }\\
\hline
3. User selects students and marks their attendance status.&4. System processes the request and sends the attendance data to the timesheet service to update records for selected students.&
\par{Request: 
        \vspace{-0.8em}
        \begin{itemize}
            \item List of selected students.
            \vspace{-0.8em}
            \item Attendance status.
            \vspace{-2.5em}
        \end{itemize}
    }\\
\hline
&5. System pushes an attendance notification to parents' phones and saves it to the database.&\\
\hline
&6. System reloads the list of student and diplay to user.&\\
\hline
\end{longtable}

\par\indent\indent\indent\parbox{0.87\textwidth}{\textbf{Mark student presence on buses}: Similar to classroom attendance, but it has several key differences. This features enables 
transportation staff record student presence at bus stops with automatic geolocation capture.}
\vspace{0.5em}
\newpage
\par\indent\indent\textbf{B. Alternative flow}
\vspace{-0.5em}
\renewcommand{\arraystretch}{1.5}
\begin{longtable}{|>{\raggedright\arraybackslash}p{4.5cm}|
                  >{\raggedright\arraybackslash}p{6cm}|
                  >{\raggedright\arraybackslash}p{4.5cm}|}                  
\caption{Mark student presence alternative flow: Mark attendance in the past.}
\label{tab:student_attendance_past_date_alternative_flow} \\
\hline
\centering\textbf{User action} & \centering\textbf{System action} & \centering\arraybackslash\textbf{Data} \\
\hline
\endfirsthead
\endhead
1. User selects a day in the past.& 2. System prevents user from selecting students to mark presence in the past.&\\
\hline
\end{longtable}
\vspace{-0.5cm}
\renewcommand{\arraystretch}{1.5}
\begin{longtable}{|>{\raggedright\arraybackslash}p{4.5cm}|
                  >{\raggedright\arraybackslash}p{6cm}|
                  >{\raggedright\arraybackslash}p{4.5cm}|}                  
\caption{Mark student presence alternative flow: Mark attendance in future date}
\label{tab:student_attendance_future_date_alternative_flow} \\
\hline
\centering\textbf{User action} & \centering\textbf{System action} & \centering\arraybackslash\textbf{Data} \\
\hline
\endfirsthead
\endhead
1. User selects the calendar.& 2. System displays a calendar with disable future date to prevent marking attendance for future date.&\\
\hline
\end{longtable}
\vspace{-0.5cm}
\renewcommand{\arraystretch}{1.5}
\begin{longtable}{|>{\raggedright\arraybackslash}p{4.5cm}|
                  >{\raggedright\arraybackslash}p{6cm}|
                  >{\raggedright\arraybackslash}p{4.5cm}|}                  
\caption{Mark student presence alternative flow: Contact parent.}
\label{tab:student_attendance_contact_parent_alternative_flow} \\
\hline
\centering\textbf{User action} & \centering\textbf{System action} & \centering\arraybackslash\textbf{Data} \\
\hline
\endfirsthead
\endhead
1. User selects the call parents option.& 2. System retrieves the parent contact information of the selected student.&\\
\hline
&3. If parent contact information exists, the system initiates a call to their mobile number, else it displays an error message.&\\
\hline
\end{longtable}

\indent\textbf{2.4.3.3 Notes}\\
\indent\indent There is no note.
\vspace{0.5em}

\indent\textbf{2.4.3.4 Pre-conditions}\\
\indent\indent The student must be either in the class or registered to a bus stop.

\vspace{0.5em}
\indent\textbf{2.4.3.5 Post-conditions}\\
\vspace{-0.7em}
\par\indent\indent\parbox{0.87\textwidth}{
    On success, the system updates student attendance records into the database and notifies parents via their mobile phones. Otherwise, it returns an error message without making any updates.
}\\

% STUDENT ATTENDANCE MANAGEMENT: GENERATE STUDENT ATTENDANCE REPORT
\subsubsection{Attendance Management: Generate Student Attendance Report}
\indent\indent\textbf{2.4.4.1 Brief description}
\vspace{0.5em}
\par\indent\indent\parbox{0.9\textwidth}{
    This use case demonstrates how administrators, authorized staff, and teachers generate a student attendance report through the system's web app, and how the system processes the request.
}\\

\vspace{0.5em}
\indent\textbf{2.4.4.2 Flow of event}
\vspace{0.5em}
\par\indent\indent\textbf{A. Basic flow}
\vspace{-0.5em}
\renewcommand{\arraystretch}{1.5}
\begin{longtable}{|>{\raggedright\arraybackslash}p{4.5cm}|
                  >{\raggedright\arraybackslash}p{6cm}|
                  >{\raggedright\arraybackslash}p{4.5cm}|}                  
\caption{View student attendance record for teacher or bus manager basic flow.}
\label{tab:student_attendance_view_attendance_basic_flow}\\
\hline
\centering\textbf{User action} & \centering\textbf{System action} & \centering\arraybackslash\textbf{Data} \\
\hline
\endfirsthead
\endhead
1. User navigates to "Student Attendance Report" function. & 2. The system retrieves the
departments which are managed by the user.& 
    \par{Response: 
        \vspace{-0.8em}
        \begin{itemize}
            \item A list of departments.
            \vspace{-1.2em}
        \end{itemize}
    }\\
\hline
3. User selects a department and select the date range.&4. System processes the request and sends the data to timesheet service.&
\par{Request: 
        \vspace{-0.8em}
        \begin{itemize}
            \item Id of the class
            \vspace{-0.8em}
            \item Selected period
            \vspace{-1.5em}
        \end{itemize}
    }\\
\hline
&5. The timesheet service fetches attendance data, then maps each student's attendance status to every date in the selected period.&
\par{Response: 
        \vspace{-0.8em}
        \begin{itemize}
            \item Student attendance records.
            \vspace{-1.5em}
        \end{itemize}
    }\\
\hline
&6. The process completes and displays the data to the user.&\\
\hline
\end{longtable}
\par\indent\indent\textbf{B. Alternative flow}
\begin{adjustwidth}{4.5em}{0em}
    Follow the Standard Alternative Flow in \textbf{\autoref{sec:common_usecase} \nameref{sec:common_usecase}}
\end{adjustwidth}
% \begin{adjustwidth}{4.5em}{0em}
%     \textbf{No attendance records found}: When a user navigates to the attendance function and no attendance records are available, the system responds differently based on the user's role. If the user is a teacher or bus manager, the system displays a list of students in their assigned class or bus route, but shows these students with no attendance status recorded. This allows educational staff to see which students are enrolled but haven't had their attendance marked yet. Alternatively, if the user is a parent accessing the system, the system displays a message box indicating that no attendance records are currently available for their child, providing clear communication about the absence of attendance data without showing unnecessary student lists.
% \end{adjustwidth}

\vspace{0.5em}
\indent\textbf{2.4.4.3 Notes}\\
\indent\indent There is no note.
\vspace{0.5em}

\indent\textbf{2.4.4.4 Pre-conditions}\\
\indent\indent The user must be assigned to at least one class.
\vspace{0.5em}

\indent\textbf{2.4.4.5 Post-conditions}\\
\vspace{-0.7em}
\par\indent\indent\parbox{0.9\textwidth}{
    If the use case is successful, student attendance records are retrieved and displayed to the user. Otherwise, the system displays an error message.
}\\
% LEAVE REQUEST MANAGEMENT: CREATE LEAVE REQUEST
\subsubsection{Leave Request Management: Create Leave Requests}
\indent\indent\textbf{2.4.5.1 Brief description}
\vspace{0.5em}
\par\indent\indent\parbox{0.9\textwidth}{
    This use case describes how leave requests can be created via mobile app (by parents/teachers) and how the system processes the submitted requests.
}\\

\vspace{0.5em}
\indent\textbf{2.4.5.2 Flow of event}
\vspace{0.5em}
\par\indent\indent\textbf{A. Basic flow}
\vspace{-0.5em}
\renewcommand{\arraystretch}{1.5}
\begin{longtable}{|>{\raggedright\arraybackslash}p{4.5cm}|
                  >{\raggedright\arraybackslash}p{6cm}|
                  >{\raggedright\arraybackslash}p{4.5cm}|}                  
\caption{Create leave request on web application.}
\label{tab:leave_request_management_create_on_web_basic_flow}\\
\hline
\centering\textbf{User action} & \centering\textbf{System action} & \centering\arraybackslash\textbf{Data} \\
\hline
\endfirsthead
\endhead

1. User opens 'Leave Request' and clicks 'Create'. & 2. The system displays create leave request screen. &\\ 
\hline
& 3. System retrieves absence types.& 
    \par{Response:
        \vspace{-0.8em}
        \begin{itemize}
            \item Absence types
            \vspace{-1.2em}
        \end{itemize}
    } \\
\hline
4. User fills all the required fields and clicks the "Save" button.&5. System processes the request and validate the leave requests (e.g., check for overlapping leave request, valid date range, maximum absence days, holiday conflicts).&
\par{Request: 
        \vspace{-0.8em}
        \begin{itemize}
            \item Leave request.
        \end{itemize}
    }\\
\hline
&6. If the leave request is valid, the system saves it to database.&\\
\hline
&7. System resfreshes the list of leave requests and displays to user.&
\par{Response: 
        \vspace{-0.8em}
        \begin{itemize}
            \item Leave requests.
        \end{itemize}
    }\\
\hline
\end{longtable}

\par\indent\indent\textbf{B. Alternative flow}
\vspace{0.3cm}
\par\indent\indent\indent\parbox{0.9\textwidth}{\textbf{Required fields missing}: When a user creates a leave request but does not provide all required fields, the system disables the Save button, preventing the user from submitting an incomplete request. It ensures all necessary information is provided and maintains data integrity.}
\vspace{0.05cm}
\par\indent\indent\indent\parbox{0.9\textwidth}{\textbf{Invalid leave request}:}
\vspace{-2em}
\renewcommand{\arraystretch}{1.5}
\begin{longtable}{|>{\raggedright\arraybackslash}p{4.5cm}|
                  >{\raggedright\arraybackslash}p{6cm}|
                  >{\raggedright\arraybackslash}p{4.5cm}|}                  
\caption{Create leave requests alternative flow: Invalid leave request.}
\label{tab:leave_request_management_invalid_alternative_flow} \\
\hline
\centering\textbf{User action} & \centering\textbf{System action} & \centering\arraybackslash\textbf{Data} \\
\hline
\endfirsthead
\endhead
5. User fills all required fields and taps "Save" button.& 6. System receives the leave request and validates it.&\\
\hline
& 6a. If validation fails (overlapping requests, exceeds remaining/maximum days, holiday conflicts, or payroll cycle issues), system returns error message.&
\par{Request: 
        \vspace{-0.8em}
        \begin{itemize}
            \item Bad request message.
        \end{itemize}
    }
\\
\hline
\end{longtable}

\indent\textbf{2.4.5.3 Notes}\\
\indent\indent There is no note.
\vspace{0.5em}

\indent\textbf{2.4.5.4 Pre-conditions}\\
\indent\indent Follow the Standard Pre-conditions in \textbf{\autoref{sec:common_usecase}}
\vspace{0.5em}

\indent\textbf{2.4.5.5 Post-conditions}\\
\vspace{-0.7em}
\par\indent\indent\parbox{0.9\textwidth}{
    If the use case is successful, the system saves the leave request to the database and notifies teachers via their mobile phones. Otherwise, the system displays an error message without creating the request.
}\\

% LEAVE REQUEST MANAGEMENT: CONFIRM LEAVE REQUEST
\subsubsection{Leave Request Management: Confirm Leave Requests}
\indent\indent\textbf{2.4.6.1 Brief description}
\vspace{0.5em}
\par\indent\indent\parbox{0.9\textwidth}{
    This use case shows teachers approve, reject, or revoke leave requests for students through the mobile applications.
}\\

\vspace{0.5em}
\indent\textbf{2.4.6.2 Flow of event}
\vspace{0.5em}
\par\indent\indent\textbf{A. Basic flow}
\vspace{-0.5em}
\renewcommand{\arraystretch}{1.5}
\begin{longtable}{|>{\raggedright\arraybackslash}p{4.5cm}|
                  >{\raggedright\arraybackslash}p{6cm}|
                  >{\raggedright\arraybackslash}p{4.5cm}|}                  
\caption{Confirm leave request on mobile application.}
\label{tab:leave_request_management_confirm_request_on_mobile_basic_flow}\\
\hline
\centering\textbf{User action} & \centering\textbf{System action} & \centering\arraybackslash\textbf{Data} \\
\hline
\endfirsthead
\endhead

1. User navigates to "Approve Leave request" function. & 2. The system retrieves a list of leave request of students which are managed by the user.&
    \par{Response:
        \vspace{-0.8em}
        \begin{itemize}
            \item List of leave request
            \vspace{-1.5em}
        \end{itemize}
    } \\
\hline
&3. System retrieves the leaves request displays to user in 3 seperated groups: pending, approved and declined.& \\
\hline
4. User approves, rejectes, or revokes a leave request.&5. System updates the leave request status.&
\par{Request: 
        \vspace{-0.8em}
        \begin{itemize}
            \item Selected leave request.
            \vspace{-1.2em}
        \end{itemize}
    }\\
\hline
&6. System resfreshes the list of leave requests and displays the updated list to user.&
\par{Response: 
        \vspace{-0.8em}
        \begin{itemize}
            \item List of updated leave request.
            \vspace{-1.2em}
        \end{itemize}
    }\\
\hline
\end{longtable}

\par\indent\indent\textbf{B. Alternative flow}
\begin{adjustwidth}{4.5em}{0em}
    Follow the Standard Alternative Flow in \textbf{\autoref{sec:common_usecase} \nameref{sec:common_usecase}}
\end{adjustwidth}
\vspace{0.5em}

\indent\textbf{2.4.6.3 Notes}\\
\indent\indent There is no note.
\vspace{0.5em}

\indent\textbf{2.4.6.4 Pre-conditions}\\
\indent\indent Follow the Standard Pre-conditions in \textbf{\autoref{sec:common_usecase}}

\vspace{0.5em}
\indent\textbf{2.4.6.5 Post-conditions}\\
\vspace{-0.7em}
\par\indent\indent\parbox{0.9\textwidth}{
    On success, the system updates the leave request status in the database and notifies parents via mobile phones (for student requests). Otherwise, the system displays an error message without any updates.
}\\


% NEWSFEED MANAGEMENT: POST NEWSFEED
\subsubsection{Newsfeed Management: Post Newsfeed}
\indent\indent\textbf{2.4.7.1 Brief description}
\vspace{0.5em}
\par\indent\indent\parbox{0.9\textwidth}{
   This use case shows how teachers create posts with images, videos, and files via the system's mobile application and how the system processes them.
}\\
\newpage
\indent\textbf{2.4.7.2 Flow of event}
\vspace{0.5em}
\par\indent\indent\textbf{A. Basic flow}
\vspace{-0.5em}
\renewcommand{\arraystretch}{1.5}
\begin{longtable}{|>{\raggedright\arraybackslash}p{4.5cm}|
                  >{\raggedright\arraybackslash}p{6cm}|
                  >{\raggedright\arraybackslash}p{4.5cm}|}                  
\caption{Post newsfeed basic flow.}
\label{tab:newsfeed_management_post_newsfeed_basic_flow}\\
\hline
\centering\textbf{User action} & \centering\textbf{System action} & \centering\arraybackslash\textbf{Data} \\
\hline
\endfirsthead
\endhead

1. User selects a class and taps "Add" button. & 2. The system displays the post creation screen.& \\
\hline
3. User writes post content, adds attachments, and taps the "Post" button. & 4. System receives the POST request from the user and validates the attachment size of the post.& 
\par{Request:
        \vspace{-0.8em}
        \begin{itemize}
            \item Post content
            \vspace{-0.8em}
            \item List of attachment in Base64 format
            \vspace{-1.5em}
        \end{itemize}
    }\\
\hline
&5. The system saves the post content to database and attachments to server.&\\
\hline
&6. System refreshes the newsfeed list (page 0, size 10).&
\par{Request: 
        \vspace{-0.8em}
        \begin{itemize}
            \item Page number, size
        \end{itemize}
    }
\vspace{-0.8em}
\par{Response: 
    \vspace{-0.8em}
    \begin{itemize}
        \item Newsfeed posts.
        \vspace{-1.5em}
    \end{itemize}
}\\
\hline
\end{longtable}

\par\indent\indent\textbf{B. Alternative flow}
\vspace{-0.5em}
\renewcommand{\arraystretch}{1.5}
\begin{longtable}{|>{\raggedright\arraybackslash}p{4.5cm}|
                  >{\raggedright\arraybackslash}p{6cm}|
                  >{\raggedright\arraybackslash}p{4.5cm}|}                  
\caption{Post newsfeed alternative flow.}
\label{tab:newsfeed_management_post_newsfeed_alternative_flow}\\
\hline
\centering\textbf{User action} & \centering\textbf{System action} & \centering\arraybackslash\textbf{Data} \\
\hline
\endfirsthead
\endhead
&5. System rejects post because the total attachments exceed 100MB, and return a bad request message.&
\par{Response:
        \vspace{-0.8em}
        \begin{itemize}
            \item Bad request message
            \vspace{-1.5em}
        \end{itemize}
    }\\
\hline
6. User removes or reduces the oversized attachment and retries.& 7. System process and validate newsfeed post (similar to steps 4 - 6 in the \textbf{Post newsfeed basic flow}).&\\
\hline
\end{longtable}

\indent\textbf{2.4.7.3 Notes}\\
\indent\indent There is no note.
\vspace{0.5em}

\indent\textbf{2.4.7.4 Pre-conditions}\\
\indent\indent The total attachment size must below or equal to 100 MB.
\vspace{0.5em}

\indent\textbf{2.4.7.5 Post-conditions}\\
\vspace{-0.7em}
\par\indent\indent\parbox{0.9\textwidth}{
    If the use case is successful, the system saves newsfeed post to database and server, making it visible in the newsfeed of the class. Otherwise, the system displays an error message and no creation is performed.
}\\

% MEDICINE REMINDER MANAGEMENT: CONFIRM MEDICINE REMINDER
\subsubsection{Medicine reminder management: Confirm medicine reminder}
\indent\indent\textbf{2.4.8.1 Brief description}
\vspace{0.5em}
\par\indent\indent\parbox{0.9\textwidth}{
    This use case describes how teachers confirm medicine reminders for students through the mobile application and how the system processes the confirmation.
}\\

\vspace{0.5em}
\indent\textbf{2.4.8.2 Flow of event}
\vspace{0.5em}
\par\indent\indent\textbf{A. Basic flow}
\vspace{-0.5em}
\renewcommand{\arraystretch}{1.5}
\begin{longtable}{|>{\raggedright\arraybackslash}p{4.5cm}|
                  >{\raggedright\arraybackslash}p{6cm}|
                  >{\raggedright\arraybackslash}p{4.5cm}|}                  
\caption{Confirm medicine reminder basic flow.}
\label{tab:medicine_reminder_management_confirm_basic_flow}\\
\hline
\centering\textbf{User action} & \centering\textbf{System action} & \centering\arraybackslash\textbf{Data} \\
\hline
\endfirsthead
\endhead

1. User selects a medicine reminder to confirm.&2. The system retrieves and displays the detailed of the selected medicine reminder.&
\par{Request:
        \vspace{-0.8em}
        \begin{itemize}
            \item Reminder ID
            \vspace{-0.8em}
        \end{itemize}
    }
    \par{Response:
        \vspace{-0.8em}
        \begin{itemize}
            \item Medicine reminder details
            \vspace{-1.5em}
        \end{itemize}
    }\\
\hline
3. User selects medicines received from parents, adds a confirmation note, and confirms the reminder.&4. System updates the medicine reminder status and saves the confirming user's information.&
\par{Request:
        \vspace{-0.8em}
        \begin{itemize}
            \item Reminder ID
            \vspace{-0.8em}
            \item Confirmation note
            \vspace{-0.8em}
            \item Received medicine items
            \vspace{-1.5em}
        \end{itemize}
    }\\
\hline
& 5. The system pushes a confirmation notification to parents.&\\
\hline
\end{longtable}

\par\indent\indent\textbf{B. Alternative flow}
\vspace{-0.5em}
\renewcommand{\arraystretch}{1.5}
\begin{longtable}{|>{\raggedright\arraybackslash}p{4.5cm}|
                  >{\raggedright\arraybackslash}p{6cm}|
                  >{\raggedright\arraybackslash}p{4.5cm}|}                  
\caption{Confirm medicine reminder alternative flow.}
\label{tab:medicine_reminder_management_alternative_flow}\\
\hline
\centering\textbf{User action} & \centering\textbf{System action} & \centering\arraybackslash\textbf{Data} \\
\hline
\endfirsthead
\endhead
1. User selects a medicine reminder to confirm.&2a. System validates the confirmation date against with the usage date of the reminder.&\\
\hline
&2b. If the confirmation date is after the medicine's usage date, the system blocks the confirmation and displays an error message.&\\
\hline
\end{longtable}

\indent\textbf{2.4.8.3 Notes}
\par\indent\indent\parbox{0.9\textwidth}{
    Once confirmed, the medicine reminder status cannot be reverted to pending.
}
\vspace{0.5em}

\indent\textbf{2.4.8.4 Pre-conditions}\\
\indent\indent  Teachers must be assigned at lease one class.
\vspace{0.25em}
\par\indent\indent\parbox{0.9\textwidth}{
    The confirmation date must not be after the medicine usage date.
}
\vspace{0.5em}

\indent\textbf{2.4.8.5 Post-conditions}\\
\vspace{-0.7em}
\par\indent\indent\parbox{0.9\textwidth}{
    If the use case is successful, the system updates the medicine reminder status, and saves information of the confirmation user. In case of failure, the system displays an error message and no updation is performed.
}\\

% NOTIFICATION MANAGEMENT: PUSH NOTIFICATION
\subsubsection{Notification management: Push notification}
\indent\indent\textbf{2.4.9.1 Brief description}
\vspace{0.5em}
\par\indent\indent\parbox{0.9\textwidth}{
    This use case describes how school administrators and teachers send notifications to parents via the web application.
}\\

\indent\textbf{2.4.9.2 Flow of event}
\vspace{0.5em}
\par\indent\indent\textbf{A. Basic flow}
\vspace{-0.5em}
\renewcommand{\arraystretch}{1.5}
\begin{longtable}{|>{\raggedright\arraybackslash}p{4.5cm}|
                  >{\raggedright\arraybackslash}p{6cm}|
                  >{\raggedright\arraybackslash}p{4.5cm}|}                  
\caption{Push notification basic flow.}
\label{tab:notification_management_push_basic_flow}\\
\hline
\centering\textbf{User action} & \centering\textbf{System action} & \centering\arraybackslash\textbf{Data} \\
\hline
\endfirsthead
\endhead
1. User navigates to the "Push notification" tab.&2. System validate user role and retrieves the list of departments which are managed by the user.&
    \par{Response:
        \vspace{-0.8em}
        \begin{itemize}
            \item List of departments
            \vspace{-1.2em}
        \end{itemize}
    }\\
\hline
3. User selects a department and specifies the date.&4a. System validates user role and retrieves all notifications for the selected department and date.&
\par{Request:
        \vspace{-0.8em}
        \begin{itemize}
            \item Id of the selected department
        \end{itemize}
    }
    \vspace{-0.8em}
    \par{Response:
        \vspace{-0.8em}
        \begin{itemize}
            \item List of notifications
            \vspace{-1.5em}
        \end{itemize}
    }\\
\hline
&4b. System validates user role and retrieves all student of the selected deparment.&
    \par{Response:
        \vspace{-0.8em}
        \begin{itemize}
            \item List of students
        \end{itemize}
    }\\
\hline
5. User selects notification to send and students to receive notifications.& 6. System receives a POST request from user and connects to notification service.&
\par{Request:
        \vspace{-0.8em}
        \begin{itemize}
            \item User session data
            \vspace{-0.8em}
            \item Notification data
        \end{itemize}
    }\\
\hline
&7. System connects to Firebase message service.&
\par{Request:
        \vspace{-0.8em}
        \begin{itemize}
            \item Authentication data
            \vspace{-1.2em}
        \end{itemize}
    }\\
\hline
&8. System finds device token of parents for each selected student, and sends notification using these device tokens.&
\par{Notification data}\\
\hline
\end{longtable}

\par\indent\indent\textbf{B. Alternative flow}
\begin{adjustwidth}{4.5em}{0em}
    Follow the Standard Alternative Flow in \textbf{\autoref{sec:common_usecase} \nameref{sec:common_usecase}}
\end{adjustwidth}
\vspace{0.5em}

\indent\textbf{2.4.9.3 Notes}\\
\indent\indent There is no note.
\vspace{0.5em}

\indent\textbf{2.4.9.4 Pre-conditions}\\
\indent\indent Teacher must be assigned at least one class.
\vspace{0.25em}
\par\indent\indent\parbox{0.9\textwidth}{
    The Firebase app must be created and integrated with the notification service.
}
\vspace{0.5em}

\indent\textbf{2.4.9.5 Post-conditions}\\
\vspace{-0.7em}
\par\indent\indent\parbox{0.9\textwidth}{
    If the use case is successful, the system pushes notifications to the parents. Otherwise, the system displays an error message.
}\\

% BUS MANAGEMENT: TRACK BUS
\subsubsection{Bus management: Track buses}
\indent\indent\textbf{2.4.10.1 Brief description}
\vspace{0.5em}
\par\indent\indent\parbox{0.9\textwidth}{
    This use case shows how parents can track the registered bus in real time through the system's mobile application.
}\\

\vspace{0.5em}
\indent\textbf{2.4.10.2 Flow of event}
\vspace{0.5em}
\par\indent\indent\textbf{A. Basic flow}
\vspace{-0.5em}
\renewcommand{\arraystretch}{1.5}
\begin{longtable}{|>{\raggedright\arraybackslash}p{4.5cm}|
                  >{\raggedright\arraybackslash}p{6cm}|
                  >{\raggedright\arraybackslash}p{4.5cm}|}                  
\caption{Track buses basic flow.}
\label{tab:bus_management_track_buses_basic_flow}\\
\hline
\centering\textbf{User action} & \centering\textbf{System action} & \centering\arraybackslash\textbf{Data} \\
\hline
\endfirsthead
\endhead
1. The user accesses the Bus function and selects the active bus route.&2a. The system sends an HTTP GET request to upgrade the connection to WebSocke.&
\par{Request:
        \vspace{-0.8em}
        \begin{itemize}
            \item GET request
            \vspace{-0.8em}
        \end{itemize}
    }\\
\hline
&2b. The system sends a STOMP CONNECT frame to establish a connection between the parent’s phone and the server.&
\par{Request:
        \vspace{-0.8em}
        \begin{itemize}
            \item CONNECT frame
        \end{itemize}
    }\\
\hline
&3a. The system sends a STOMP SEND frame to the server to identify the transportation staff managing the selected bus route.&
    \par{Request:
        \vspace{-0.8em}
        \begin{itemize}
            \item Bus route id
        \end{itemize}
    }\\
\hline
& 3b. The system sends a STOMP SUBSCRIBE frame to receive future location updates from the transportation staff.&Same as above\\
\hline
& 3c. Once the transportation staff is identified, the system emits a STOMP MESSAGE frame to all subscribers.&Same as above\\
\hline
&4a. When the transportation staff receives the message, the system retrieves their current location (longitude and latitude) every 50 meters and sends a SEND frame to the server.&
\par{Request:
        \vspace{-0.8em}
        \begin{itemize}
            \item Longitude and latitude
            \vspace{-1.2em}
        \end{itemize}
    }\\
\hline
&4b. The system emits a MESSAGE frame to the topic previously subscribed to by the parent’s phone.&
\par{Reponse:
        \vspace{-0.8em}
        \begin{itemize}
            \item Longitude and latitude
            \vspace{-1.2em}
        \end{itemize}
    }\\
\hline
&5. The system receives the location data and displays it to the user.&\\
\hline
\end{longtable}

\par\indent\indent\textbf{B. Alternative flow}
\vspace{0.6em}
\par\indent\indent\indent\parbox{0.9\textwidth}{\textbf{Transportation Staff Not Found}: When a user wants to track the bus location that transports their child but the system cannot identify any transportation staff assigned to that route, it will display an error message to notify the user and ensure transparency.}
\vspace{0.5em}
\par\indent\indent\indent\parbox{0.9\textwidth}{\textbf{GPS Not Enabled}:  If the transportation does not enable the mobile device GPS, the system cannot get their current location. In this case, the system will display an error message: Cannot get the location of transportation staff.}
\vspace{0.1em}

\indent\textbf{2.4.10.3 Notes}\\
\indent\indent There is no note.
\vspace{0.5em}

\indent\textbf{2.4.10.4 Pre-conditions}\\
\indent\indent Children must be registered to a bus route.
\vspace{0.25em}
\par\indent\indent\parbox{0.9\textwidth}{
    The transportation staff must let the app in the foreground or background state and enable the GPS.
}

\vspace{1em}
\indent\textbf{2.4.10.5 Post-conditions}\\
\vspace{-0.7em}
\par\indent\indent\parbox{0.9\textwidth}{
    If the use case is successful, parents can track the buses in real-time. However, if the use case fails, the system will display to users, and they cannot see the location of the bus in real-time.
}\\


% EMPLOYEE ATTENDANCE
% \subsubsection{Employee attendance management: Export employee attendance data}
% \indent\indent\textbf{2.5.8.1 Brief description}
% \vspace{0.5em}
% \par\indent\indent\parbox{0.9\textwidth}{
%     The "Export employee attendance data" is a feature available in web browsers that allows school administrators, staff and teacher to export detailed check-in/check-out records of employees under their supervision in Excel. The exported file contains daily check-in/check-out times, total working hours, and absence records. This feature enables user to export attendance data in a specific month.
% }\\

% \vspace{0.5em}
% \indent\textbf{2.5.8.2 Flow of event}
% \vspace{0.5em}
% \par\indent\indent\textbf{A. Basic flow}
% \renewcommand{\arraystretch}{1.5}
% \begin{longtable}{|>{\raggedright\arraybackslash}p{4.5cm}|
%                   >{\raggedright\arraybackslash}p{6cm}|
%                   >{\raggedright\arraybackslash}p{4.5cm}|}                  
% \caption{Export employee attendance basic flow.}
% \label{tab:employee_attendance_management_export_attendance_data_basic_flow} \\
% \hline
% \centering\textbf{User action} & \centering\textbf{System action} & \centering\arraybackslash\textbf{Data} \\
% \hline
% \endfirsthead
% \endhead

% 1. User navigates to Attendance report tab. & 2. The system gets the organizations which are managed by the user (displayed in a tree structures).& 
%     Request:
%         \vspace{-0.8em}
%         \begin{itemize}
%             \item User session data
%         \end{itemize}
%     \\
% \hline
% 3. User selects month, year and an organization from the loaded organizations to export data and clicks Export excel button.&4. System verifies the user role and get all the active and inactive employees under user supervisor from the selected organization and the child organizations.&
% \par{Request:
%         \vspace{-0.8em}
%         \begin{itemize}
%             \item User session data
%             \vspace{-0.8em}
%             \item Id of the selected organization
%             \vspace{-0.8em}
%             \item Selected month
%             \vspace{-0.8em}
%             \item Selected year
%         \end{itemize}
%     } 
%     \\
% \hline
% & 5. System sends a request to the timesheet service to get the attendance records.&
% \par{Request:
%         \vspace{-0.8em}
%         \begin{itemize}
%             \item User session data
%             \vspace{-0.8em}
%             \item List of employees
%             \vspace{-0.8em}
%             \item Selected month
%             \vspace{-0.8em}
%             \item Selected year
%         \end{itemize}
%     }\\
% \hline
% &6. The timesheet service processes the request and returns all the attendance data by the list of employee and the selectd month and year.&
% \par{Response:
%         \vspace{-0.8em}
%         \begin{itemize}
%             \item List of attendance data.
%         \end{itemize}
% }\\
% \hline
% &7. System process the list of employees attedance data.&\\
% \hline
% &8. System creates a new excel file based on the excel template from server and add the value of the attendance list to excel file.&\\
% \hline
% &9. System converts the excel file to byte array, deletes the temporary excel file just created, and return the byte array to user.&
% \par{Response:
%         \vspace{-0.8em}
%         \begin{itemize}
%             \item Byte array of the excel file.
%         \end{itemize}
% }\\
% \hline
% 10. User receive the Excel file containing attendance data of the selected month.&&\\
% \hline
% \end{longtable}

% \par\indent\indent\textbf{B. Alternative flow}
% \vspace{1em}
% \par\indent\indent\indent\parbox{0.9\textwidth}{- \textbf{No organizations available}: If no organization is available, it means the user does not manage any organization or employees. In this case, they can only export their own attendance data.}
% \renewcommand{\arraystretch}{1.5}
% \begin{longtable}{|>{\raggedright\arraybackslash}p{4.5cm}|
%                   >{\raggedright\arraybackslash}p{6cm}|
%                   >{\raggedright\arraybackslash}p{4.5cm}|}                  
% \caption{Export employee attedance data alternative flow: No orginaizations is available.}
% \label{tab:employee_attendance_management_export_attendance_no_org_alternative_flow} \\
% \hline
% \centering\textbf{User action} & \centering\textbf{System action} & \centering\arraybackslash\textbf{Data} \\
% \hline
% \endfirsthead
% \endhead
% 1. User navigates to Attendance report tab.& 2. System find no organization managed by the user and fix the user code to the employee code and user cannot delete it.&
% Request:
%         \vspace{-0.8em}
%         \begin{itemize}
%             \item User session data
%         \end{itemize}
%     \\
% \hline
% 3. User selects month, year and an organization from the loaded organizations to export data and clicks Expor excel button.&4. System find the user information in database.&
% \par{Request:
%         \vspace{-0.8em}
%         \begin{itemize}
%             \item User session data
%             \vspace{-0.8em}
%             \item Employee code of the user
%             \vspace{-0.8em}
%             \item Selected month
%             \vspace{-0.8em}
%             \item Selected year
%         \end{itemize}
%     } 
%     \\
% \hline
% &5-10. Same as processing as basic flow.&\\
% \hline
% \end{longtable}

% \par\indent\indent\indent\parbox{0.9\textwidth}{- \textbf{Error during data retrieval}:}
% \vspace{-1.5em}
% \renewcommand{\arraystretch}{1.5}
% \begin{longtable}{|>{\raggedright\arraybackslash}p{4.5cm}|
%                   >{\raggedright\arraybackslash}p{6cm}|
%                   >{\raggedright\arraybackslash}p{4.5cm}|}                  
% \caption{Export employee attedance data alternative flow: Error during data retrieval.}
% \label{tab:employee_attendance_management_export_attendance_error_retrieves_alternative_flow} \\
% \hline
% \centering\textbf{User action} & \centering\textbf{System action} & \centering\arraybackslash\textbf{Data} \\
% \hline
% \endfirsthead
% \endhead
% & 6. Timesheet service encounters an error (database connection failure, timeout, etc.).&
% Response:
%         \vspace{-0.8em}
%         \begin{itemize}
%             \item Bad request message
%         \end{itemize}
%     \\
% \hline
% &7. System display a error message box.&\\
% \hline
% \end{longtable}
% \indent\textbf{2.5.8.3 Special Requirements}\\
% \indent\indent This feature is only available on the web application of the system.
% \vspace{0.5em}

% \indent\textbf{2.5.8.4 Pre-conditions}\\
% \indent\indent User must be logged in the system with valid credentials.
% \vspace{0.25em}
% \par\indent\indent\parbox{0.9\textwidth}{
%     User must have permission to export employee data (school administrator, staff and teacher).
% }
% \vspace{0.25em}
% \par\indent\indent\parbox{0.9\textwidth}{
%     The system has access to both the database and timesheet service.
% }
% \vspace{0.25em}
% \par\indent\indent\parbox{0.9\textwidth}{
%     The Excel template file must be existed on the server.
% }
% \vspace{0.5em}

% \indent\textbf{2.5.8.5 Post-conditions}\\
% \vspace{-0.7em}
% \par\indent\indent\parbox{0.9\textwidth}{
%     If the use case is successful, a excel file containing employee attendance data is generated and downloaded to the user's device. Otherwise, the system returns an error message.
% }\\